\newpage
\section{Ponto de projeto e seção de referência}

O ponto de projeto escolhido para a análise e otimização da asa é a condição
de cruzeiro da aeronave, visto que esta é a situação em que a aeronave ficará
durante maior tempo dentro de sua missão típica, especificada pelos
pré-requisitos. Que são: altitude de $h = 35\,000$ ft ($10\,668$ m) operando a
um Mach de $M_\infty = 0{,}85$. Nesta condição, com os cálculos de gasto de
combustível na missão, o peso médio considerado foi de $W = 229\,669{,}3$ kgf.

Como seção de referência adota-se a estação da corda média aerodinâmica
(MAC), em $\eta = 0{,}398$, com $c_{\mathrm{ref}} = 6{,}899$ m.

\subsection*{Obtenção de $c_{\ell,\mathrm{ref}}$, $M_n$ e $Re_{\mathrm{ref}}$}

Do modelo de atmosfera ISA, temos, para a dada altitude,
$\rho_\infty = 0{,}38046 \ \mathrm{kg \cdot m^{-3}}$,
$a_\infty = 296{,}587 \ \mathrm{m \cdot s^{-1}}$ e
$\mu_\infty = 1{,}4451 \times 10^{-5} \ \mathrm{Pa \cdot s}$. Com isso:

\begin{equation*}
    C_L = \frac{W}{\tfrac{1}{2}\,\rho_\infty\,(M_\infty a_\infty)^2\,S_{\mathrm{ref}}}
        = 0{,}5053.
\end{equation*}

O $C_L$ da aeronave, porém, não é o $c_\ell$ que a seção de referência
precisa sustentar. Duas correções se aplicam. Primeiro, assumindo a
distribuição de sustentação elíptica --- o alvo de projeto da asa ---, o
coeficiente local na estação da MAC vale
\begin{equation*}
    c_\ell(\eta) = \frac{4\,C_L\,S_{\mathrm{ref}}}{\pi\,b\;c(\eta)}\sqrt{1-\eta^2}
    \quad\Rightarrow\quad c_\ell(0{,}398) = 0{,}5246,
\end{equation*}
superior ao $C_L$ porque, na asa afilada ($\lambda = 0{,}24$), a corda cai
mais depressa que a carga elíptica.

Segundo, o perfil é projetado no plano \emph{normal} à linha de enflechamento
--- é nesse plano que a onda de choque se forma. Pela teoria de asas
enflechadas, com $\Lambda_{c/2} = 32{,}158^\circ$:
\begin{equation*}
    M_n = M_\infty \cos\Lambda_{c/2} = 0{,}7196,
    \qquad
    c_{\ell,\mathrm{ref}} = \frac{c_\ell}{\cos^2\Lambda_{c/2}} = 0{,}7319.
\end{equation*}

A corda nesse plano é $c_n = c_{\mathrm{ref}}\cos\Lambda_{c/2} = 5{,}841$ m e,
com a velocidade normal $V_n = M_n\,a_\infty$, o Reynolds da seção fica
\begin{equation*}
    Re_{\mathrm{ref}} = \frac{\rho_\infty\,V_n\,c_n}{\mu_\infty}
        = 3{,}28 \times 10^{7}.
\end{equation*}

A espessura relativa da asa na estação da MAC, definida no Lab 02, é de 15\%
na direção da corrente; no plano normal,
$(t/c)_{\mathrm{ref}} = 0{,}15/\cos\Lambda_{c/2} = 0{,}1772$, que é a
espessura mínima imposta como restrição na otimização.

\begin{table}[H]
\centering
\caption{Dados para a seção de referência.}\label{tab:ref_section}
\renewcommand{\arraystretch}{1.25}
\begin{tabular}{lcl}
\toprule
Parâmetro & Valor & Descrição \\
\midrule
$h$                     & 35.000 & Altitude do ponto de projeto [ft] \\
$M_\infty$              & 0,85 & Mach do ponto de projeto \\
$W$                     & 229.669,3 & Peso da aeronave no ponto de projeto [kgf] \\
$S_{\mathrm{ref}}$      & 368,833 & Área de referência da aeronave [m$^2$] \\
$\rho_\infty$           & 0,38046 & Densidade do ar no ponto de projeto [kg/m$^3$] \\
$a_\infty$              & 296,587 & Velocidade do som no ponto de projeto [m/s] \\
$C_L$                   & 0,5053 & Coeficiente de sustentação da aeronave \\
$\Lambda_{c/2}$         & 32,158 & Enflechamento a meia corda [graus] \\
$M_n$                   & 0,7196 & Mach normal ao enflechamento \\
$c_{\mathrm{ref}}$      & 6,899 & Corda da seção de referência (MAC) [m] \\
$c_n$                   & 5,841 & Corda no plano normal [m] \\
$c_{\ell,\mathrm{ref}}$ & 0,7319 & Coef.\ de sustentação de projeto da seção (plano normal) \\
$Re_{\mathrm{ref}}$     & $3{,}28\times 10^{7}$ & Reynolds da seção de referência \\
$(t/c)_{\mathrm{ref}}$  & 0,1772 & Espessura relativa mínima (plano normal) \\
\bottomrule
\end{tabular}
\end{table}
